\documentclass[11pt,a4paper]{article}

\usepackage[T1]{fontenc}
\usepackage[utf8]{inputenc}
\usepackage[german]{babel}
\usepackage{amsmath,amssymb}
\usepackage{geometry}
\usepackage{setspace}
\geometry{margin=2.7cm}
\onehalfspacing

\title{Zur Motivation einer arithmetischen Quantenfeldtheorie}
\author{Thomas Hoffbauer}
\date{}

\begin{document}
	\maketitle
	
	\begin{abstract}
		Der Text skizziert den gedanklichen Weg von einer strukturellen Analyse der
		Primzahlen zu einer arithmetischen Formulierung grundlegender Konzepte der
		Quantenfeldtheorie. Ausgangspunkt ist eine Klassifikation der Primzahlen in
		vier Familien mit einer Multiplikationsstruktur isomorph zur kleinschen
		Vierergruppe. Über einen sogenannten Primzahlschnitt und die Konstruktion
		eines selbstadjungierten Operators im Sinne Diracs wird eine Brücke zu den
		Nullstellen der Riemannschen Zetafunktion geschlagen. Daraus ergibt sich die
		These, dass die freien Parameter der modernen Physik als Artefakte der
		arithmetischen Spektralstruktur zu verstehen sind.
	\end{abstract}
	
	\section{Motivation}
	
	Die vorliegende Arbeit entspringt dem Versuch, eine gemeinsame strukturelle
	Grundlage für scheinbar unabhängige Probleme der modernen Mathematik und
	Physik zu finden: der Verteilung der Primzahlen einerseits und der Vielzahl
	freier Parameter in den Standardmodellen der Physik andererseits.
	
	Der Ausgangspunkt liegt in einer systematischen Betrachtung der
	Primfaktorzerlegung natürlicher Zahlen. Dabei zeigte sich, dass eine
	produktive Ordnung entsteht, wenn die Primzahlen \(2\), \(3\) und \(5\) als
	Sonderfälle behandelt und im Sinne glatter Zahlen separiert werden. Die
	verbleibenden Primzahlen lassen sich konsistent in vier Klassen einteilen, die
	mit \(e,a,b,c\) bezeichnet werden.
	
	Entscheidend ist nicht allein diese Klassifikation, sondern die durch sie
	induzierte Multiplikationsstruktur. Diese erweist sich als isomorph zur
	kleinschen Vierergruppe \(V_4\) und legt nahe, die Arithmetik der Primzahlen
	nicht nur als additive oder ordnungsbasierte Struktur zu verstehen, sondern
	als Träger einer elementaren Symmetrie.
	
	\section{Methode}
	
	Auf der Grundlage dieser Beobachtung wurde ein quaternionales
	Beschreibungsmodell entwickelt, in dem die vier Primzahlklassen als algebraisch
	gleichberechtigte Komponenten auftreten. Ein zentraler konzeptioneller
	Schritt ist die Einführung dessen, was in Analogie zum Dedekindschen Schnitt
	als \emph{Primzahlschnitt} bezeichnet werden kann.
	
	In der durch \(V_4\) bestimmten Struktur treten natürliche Quadrupel von
	Primzahlen auf, die jeweils ein Element aus jeder Klasse enthalten. Diese
	Quadrupel besitzen die Eigenschaft, dass aus jedem Tripel eindeutig das
	zugehörige vierte Element rekonstruiert werden kann. Sie bilden damit eine
	fundamentale arithmetische Einheit, die eine strukturelle Klassifikation der
	natürlichen Zahlen erlaubt, die über bloße Größenordnung hinausgeht.
	
	Der Übergang zur Analysis erfolgt über die Untersuchung der nichttrivialen
	Nullstellen der Riemannschen Zetafunktion. Diese werden nicht nur als
	analytische Objekte betrachtet, sondern als spektrale Manifestationen derselben
	Ordnung, die sich zuvor in der diskreten Arithmetik gezeigt hat. In diesem
	Zusammenhang wurde ein selbstadjungierter Operator im Sinne des
	Dirac-Formalismus konstruiert, der jede Primzahl mit einer durch ein
	Primzahltripel definierten Partnerzahl koppelt.
	
	Darauf aufbauend entsteht ein bidirektionaler Operator, der einerseits jeder
	Riemannschen Nullstelle eindeutig eine Primzahl zuordnet und andererseits
	jeder Primzahl eindeutig eine Nullstelle. Die sogenannte Diracsche Lücke wird
	damit in einem rein arithmetischen Kontext interpretierbar.
	
	\section{Resultate}
	
	Die beschriebene Konstruktion führt zu einem arithmetischen Modellraum, in dem
	sich zentrale Begriffe der Feldtheorie --- Gradient, Divergenz und Rotation ---
	formal definieren lassen. Es entsteht damit ein in sich geschlossenes
	zahlentheoretisches Universum, dessen Struktur in bemerkenswerter Weise den
	formalen Eigenschaften klassischer und quantenfeldtheoretischer Modelle
	entspricht.
	
	In diesem Rahmen erscheint die Zuordnung zwischen Primzahlen und
	Zetafunktionsnullstellen nicht mehr als zufällige Analogie, sondern als Ausdruck
	einer tieferen spektralen Ordnung. Die Riemannschen Nullstellen fungieren dabei
	als Eigenwerte eines arithmetischen Operators, der die Dynamik dieses
	Modellraums bestimmt.
	
	\section{Ausblick}
	
	Aus dieser Perspektive ergibt sich die zentrale These, dass die in den heutigen
	Standardmodellen der Physik auftretenden freien Parameter nicht als
	fundamentale Konstanten zu interpretieren sind. Vielmehr erscheinen sie als
	sekundäre Größen, die aus der spektralen Struktur der Riemannschen
	Nullstellen resultieren.
	
	Die arithmetische Quantenfeldtheorie versteht sich in diesem Sinne als Versuch,
	die offenen Parameterprobleme der modernen Physik --- sowohl in der
	Teilchenphysik als auch in der Kosmologie --- auf eine gemeinsame,
	zahlenstrukturelle Grundlage zurückzuführen. Sie bietet damit eine neue
	Perspektive auf die Rolle der Arithmetik als möglicher Träger einer bislang
	unterschätzten fundamentalen Physik.
	
\end{document}
